% !PRINTED 259
\weberpage{259}

\begin{center}
\S.\,67.

\gesperrt{Congruenzkörper zweiten Grades.}
\end{center}

In jedem Congruenzkörper $\mathfrak{E}_{n,p}$ für den Modul $p$ vom $n^{\mathrm{ten}}$
Grade ist der Congruenzkörper ersten Grades $\mathfrak{E}_{1,p}$, der aus den
nach dem Modul $p$ genommenen rationalen Zahlen besteht, als
Theiler enthalten. Daher ist auch in der Gruppe der linearen
Substitutionen $E_{n,p}$ eine Gruppe als Theiler enthalten, die aus
allen Substitutionen der Form

\weq{1}{A = \begin{pmatrix} a, & b \\ c, & d \end{pmatrix}, \quad ab-cd=1 % [sic] printed ab-cd, not ad-bc
}
besteht, worin $a, b, c, d$ nach dem Modul $p$ genommene ganze
rationale Zahlen sind.

Diese Gruppe, die als die Gruppe der Modulargleichungen
in der Theorie der elliptischen Functionen auftritt$^{1)}$, ist das
eigentliche Ziel unserer Betrachtungen. Es bietet aber für die
Untersuchung dieser Gruppe, und besonders für die Aufsuchung
ihrer Theiler, den grössten Vortheil, diese Gruppe nicht selbstständig
für sich zu betrachten, sondern als Theiler einer
Gruppe $E_{2,p}$. Bezeichnen wir nämlich mit $N$ irgend einen feststehenden
quadratischen Nichtrest von $p$, so ist, wie schon oben
bemerkt, $t^2-N$ nach dem Modul $p$ irreducibel, und wir erhalten
einen Congruenzkörper $2^{\mathrm{ten}}$ Grades, wenn wir eine
\gesperrt{Galois'sche} imaginäre Zahl durch die Gleichung

\weq{2}{\varepsilon^2 = N}
einführen.

Es soll der Deutlichkeit wegen für die nächsten Betrachtungen
festgesetzt sein, dass die kleinen lateinischen Buchstaben
ganze rationale Zahlen (nach dem Modul $p$ genommen), die wir
hier auch \gesperrt{reelle Zahlen} nennen, die griechischen Buchstaben
Zahlen des Körpers $\mathfrak{E}_{2,p}$ bedeuten sollen.

Dieser Körper $\mathfrak{E}_{2,p}$ hat die für uns sehr wichtige Eigenschaft,
dass \gesperrt{in ihm alle reellen Zahlen Quadrate sind}.

Wenn nämlich $a$ quadratischer Rest von $p$ ist, so können
wir $a=x^2$ setzen, und wenn $b$ quadratischer Nichtrest ist, so
kann $x^2=bN$ befriedigt werden, also $b=(x\varepsilon^{-1})^2$.

\rule{4cm}{0.4pt}

$^{1)}$ Man vergleiche des Verfassers Buch: \glqq Elliptische Functionen und
algebraische Zahlen\grqq{}. Braunschweig 1891.
